\documentclass[11pt]{article}

% --- Geometry, fonts, micro-typography ----------------------------------------
\usepackage[margin=1in]{geometry}
\usepackage[T1]{fontenc}
\usepackage{lmodern}
\usepackage{microtype}

% --- Math, graphics, tables ---------------------------------------------------
\usepackage{amsmath,amssymb,amsthm}
\usepackage{graphicx}
\usepackage{booktabs}
\usepackage{multirow}
\usepackage{array}
\usepackage{xcolor}
\usepackage{enumitem}
\usepackage{caption}
\usepackage{subcaption}
\usepackage{tikz}
\usetikzlibrary{positioning,arrows.meta,fit,shapes.geometric,calc}

% --- Bibliography (numeric, sorted as cited) ----------------------------------
\usepackage[numbers,sort&compress]{natbib}
\usepackage{hyperref}
\hypersetup{
    colorlinks=true,
    linkcolor=blue!50!black,
    citecolor=blue!50!black,
    urlcolor=blue!50!black,
}

% --- Custom macros ------------------------------------------------------------
\newcommand{\todo}[1]{\textcolor{red}{\textbf{[TODO: #1]}}}
\newcommand{\pending}[1]{\textcolor{orange!80!black}{\textbf{#1}}}
\newcommand{\norm}[1]{\left\lVert#1\right\rVert}
\newcommand{\eqdef}{\overset{\Delta}{=}}
\newcommand{\R}{\mathbb{R}}
\newcommand{\E}{\mathbb{E}}
\newcommand{\KL}{D_{\mathrm{KL}}}

% --- Title / Author -----------------------------------------------------------
\title{\textbf{PLA: Peripersonal Language-Action Policies via\\Whole-Body Time-of-Flight Proximity Sensing}}
\author{%
  Anonymous Submission\\
  \textit{Submitted to CoRL 2026 (double-blind)}%
}
\date{}

\begin{document}
\maketitle

% =============================================================================
\begin{abstract}
\noindent
Vision-language-action (VLA) policies achieve impressive open-vocabulary
manipulation but struggle in the \emph{peripersonal} regime --- the volume
of space within tens of centimetres of the robot body, where small
perception errors translate to grasp failure or collision. We argue that
this regime is fundamentally under-served by RGB-only perception: occlusion,
foreshortening, and texture-poor surfaces are most damaging precisely where
millimetre-scale precision matters most. We present \textbf{PLA
(Peripersonal Language-Action)}, a policy architecture that augments a
frozen vision-language backbone with a learned representation of whole-body
time-of-flight (ToF) proximity sensing. PLA encodes 32 VL53L5CX-class
8$\times$8 SPAD multi-zone sensors mounted across the FR3 robot's links and
end-effector through a shared, layer-normed MLP that produces one token per
sensor. These tokens are concatenated with frozen Molmo2-4B visual-language
tokens and a proprioceptive token, then passed to an Action Chunking
Transformer (ACT) decoder that predicts a 100-step joint-delta horizon in a
single forward pass. PLA differs from a comparable ACT baseline by
\emph{exactly one} architectural switch: the presence or absence of the
proximity stream. We pre-register a 100-episode-per-condition evaluation on
four MolmoSpaces tasks (open-workspace pick-and-place, near-contact pick-and-place,
language-conditioned colour selection, and spatial-relation placement) with
paired-bootstrap p-values under a fixed seed schedule. We report Day-1
system validation (forward-pass, gradient-flow, dataset-pipeline, and
statistical-method correctness) and lock the experimental protocol prior to
data collection. \emph{Final empirical results will be reported in
Section~\ref{sec:results} on Day~14 of execution; the methods,
hyperparameters, statistical procedure, and acceptance criteria are
finalised here.}
\end{abstract}

% =============================================================================
\section{Introduction}
\label{sec:intro}

Recent vision-language-action policies~\citep{brohan2023rt2,kim2024openvla,
black2024pi0,molmo2024} have established the modern recipe for generalist
manipulation: a large pre-trained vision-language model (VLM) provides
open-vocabulary perception and language grounding, and a lightweight action
head converts contextual representations into low-level joint commands.
This recipe is correct on average. It is not yet adequate at the boundaries.

\paragraph{The peripersonal regime.}
Cognitive scientists use the term \emph{peripersonal space} to denote the
volume immediately surrounding the body, monitored by multimodal neurons in
parietal cortex specifically tuned for action and defense
\citep{rizzolatti1981peripersonal,graziano2009peripersonal}. In humans this
volume is rich in tactile and proximity cues; in robots it is the regime in
which (i) RGB foreshortening is most severe, (ii) occlusion is unavoidable
because the robot itself blocks the camera, and (iii) error tolerances are
millimetres rather than centimetres. The same VLA that can correctly point
to ``the red mug'' from a metre away may fail to insert its gripper into
the mug's handle from 4\,cm.

\paragraph{Why this is a representation problem, not a control problem.}
A common rebuttal is that fine motion is the controller's responsibility,
not the policy's. Modern force-controlled arms can compensate for small
sensing errors at runtime. But this assumes the policy has \emph{committed}
to an action sequence that is feasible given the local geometry. In the
peripersonal regime, infeasibility manifests not as a small tracking error
but as collision or grasp miss --- a categorical failure that no controller
can rescue. The policy needs a representation of geometry that does not
share the failure modes of its primary visual stream.

\paragraph{Whole-body time-of-flight.}
We consider \emph{multi-zone time-of-flight (ToF) proximity sensors}
distributed over the robot body. The VL53L5CX class of sensors used in
GenTact-style robotic skins~\citep{patel2023gentact} produces an
8$\times$8 grid of SPAD-based depth measurements at 45$^\circ$ FOV and
20--4000\,mm range. Mounted across multiple links, dozens of sensors
provide a spatially distributed view of nearby geometry that is
\emph{complementary} to RGB: ToF is unaffected by texture, lighting, or
the foreshortening near the gripper, but is poor at object identity and
language grounding. We argue this is exactly the right complement to a
VLM-backed policy.

\paragraph{Contributions.}
\begin{enumerate}[leftmargin=1.5em,topsep=0.2em,itemsep=0.2em]
\item \textbf{Architecture (Section~\ref{sec:method}).} We propose PLA, a
  policy that fuses frozen Molmo2-4B visual-language tokens with learned
  whole-body ToF proximity tokens via concatenation followed by a shared
  ACT decoder. The proximity encoder is a single shared MLP over
  flattened 8$\times$8 zone readings; per-sensor identity is preserved by
  the per-token slot in the decoder context.
\item \textbf{The one-flag baseline (Section~\ref{sec:method-baseline}).}
  PLA and the VLM-only ACT baseline differ only in the presence of the
  proximity stream. They share architecture, hyperparameters, dataset,
  and seed schedule. This makes the proximity contribution
  \emph{causally identifiable} from the comparison.
\item \textbf{Sim-to-data pipeline (Section~\ref{sec:system-data}).} We
  detail the full pipeline from URDF/MJCF skin construction through
  per-step camera-based depth rendering with a calibrated noise model, to
  HDF5 trajectory storage with schema validation and a
  proximity-informative coverage check that gates training start.
\item \textbf{Pre-registered evaluation (Section~\ref{sec:protocol}).} We
  finalise tasks, sample sizes, statistical methodology
  (paired-bootstrap with shared seeds, $\alpha = 0.05$, two-sided), and
  acceptance criteria \emph{before} running the headline experiment.
\item \textbf{Open implementation.} An importable Python package,
  configurable from flat YAML files with one switch separating PLA from
  baseline, supports collection, training, evaluation, and ablation
  through a single training entry point.
\end{enumerate}

% =============================================================================
\section{Related Work}
\label{sec:related}

\paragraph{Vision-language-action policies.}
RT-2~\citep{brohan2023rt2} cast manipulation as autoregressive token
prediction over a co-fine-tuned VLM. OpenVLA~\citep{kim2024openvla}
provided an open recipe for VLA training. $\pi_0$~\citep{black2024pi0}
combined a VLM with a flow-matching action head; Molmo
\citep{molmo2024} added pixel-level pointing supervision that has proved
useful for spatial grounding. None of these works model proximity as a
first-class input. Our contribution is orthogonal: PLA composes with any
of these backbones (we use Molmo2 for its strong pointing prior), and
adds a stream that is decoupled from the visual one.

\paragraph{Action chunking and chunked decoders.}
ACT~\citep{zhao2023learning} predicts an action chunk in one forward
pass via a CVAE encoder over the future chunk plus a transformer
decoder. Diffusion Policy~\citep{chi2023diffusion} replaces the CVAE
with iterative denoising. We use ACT because (i) one-shot inference fits
within our 5\,Hz control budget without distillation, and (ii) the CVAE
bottleneck regularises against multimodal expert demos arising from
heuristic TAMP planners (Section~\ref{sec:design-act}).

\paragraph{Tactile and proximity sensing.}
GelSight~\citep{yuan2017gelsight} and DIGIT~\citep{lambeta2020digit}
provide high-resolution contact information at fingertips. ReSkin
\citep{bhirangi2021reskin} extends this to whole-hand tactile sheets.
GenTact~\citep{patel2023gentact} introduced a parametric Blender
workflow for designing whole-body ToF skins on arbitrary robot meshes.
None of these works couple proximity to a VLM backbone for
language-conditioned manipulation. Concurrent work explores tactile-VLA
\citep{tactilevla2025} but with a single-fingertip sensor; we operate at
whole-body scale.

\paragraph{Peripersonal space in robotics.}
Roncone et al.~\citep{roncone2016peripersonal} formalised peripersonal
maps via touch-evoked optical-flow associations on the iCub. Our approach
inherits the conceptual claim --- a space-around-the-body that is
behaviourally and representationally distinct from extrapersonal vision
--- and supplies the missing piece: a \emph{learnable} encoding compatible
with modern VLA stacks.

% =============================================================================
\section{Background and Motivation}
\label{sec:background}

\subsection{Failure modes of RGB-only VLAs in the peripersonal regime}
\label{sec:bg-failure}

We characterise three failure modes that motivate the proximity stream:

\begin{enumerate}[leftmargin=1.5em,topsep=0.2em,itemsep=0.2em]
\item \textbf{Self-occlusion.} Once the gripper enters the camera frame's
  near field, the arm itself occludes the target. The VLM has to predict
  finger position from a partial view of the gripper.
\item \textbf{Texture poverty.} Many household objects are uniformly
  coloured (mugs, cups, fruits). At the close range needed for grasp
  alignment, surface texture cues vanish, leaving the VLM with only
  silhouette information.
\item \textbf{Cluttered free-space estimation.} Standard RGB depth
  estimators (MiDaS, ZoeDepth) are trained on scene-scale data and have
  high error at the 1--10\,cm range relevant for collision avoidance.
\end{enumerate}

ToF proximity sensing is, by construction, immune to all three: it
measures geometry directly at exactly the range where the failure modes
bite hardest, and the per-sensor field of view samples \emph{around} the
arm rather than \emph{through} it.

\subsection{What we expect proximity to help with --- and what it cannot}

We are explicit about the scope of the contribution. Proximity should
help wherever local geometry, not object identity, is the binding
constraint: near-contact obstacle avoidance, gripper-target alignment,
and place-feasibility checks. It should \emph{not} help with
language-conditioned object selection or spatial relations between
distant objects. Our task suite (Section~\ref{sec:protocol-tasks}) tests
both --- we expect a large gap on near-contact and a small or zero gap
on language-only tasks. A method that wins everywhere is suspicious;
proximity tokens that ``help'' on \texttt{pnp\_color} would suggest the
encoder is leaking task-irrelevant signal (e.g.\ encoding object
identity from reflectance) rather than measuring proximity.

% =============================================================================
\section{System}
\label{sec:system}

\subsection{Sensor skin}
\label{sec:system-skin}

Our skin design comprises 32 VL53L5CX-class sensors distributed across
the Franka FR3's links 2, 3, 5, 6 and end-effector, with denser coverage
on link 6 and the gripper (8 sensors) where the geometry-informativeness
density is highest. Each sensor is mounted with its optical axis along
the body-frame $+Z$ direction (outward) on a 3\,mm offset stand-off to
clear the substrate. The skin is designed parametrically in
Blender~\citep{blender} using the GenTact toolbox~\citep{patel2023gentact},
exported as a JSON of $\{$name, position, Euler$\}$ tuples, and
materialised in MuJoCo~\citep{todorov2012mujoco} as one fixed
\verb|<camera>| body per sensor, set to 8$\times$8 resolution and
45$^\circ$ vertical FOV.

\subsection{Sensor model and rendering}
\label{sec:system-model}

At each simulation step, the \verb|ToFSensorArray| module renders all
sensor cameras using MuJoCo's depth renderer. The procedure is:

\begin{enumerate}[leftmargin=1.5em,topsep=0.2em,itemsep=0.2em]
\item For each sensor camera (in MJCF enumeration order, \emph{not}
  lexicographic), render an 8$\times$8 depth image in metres.
\item Convert to millimetres and clip to $[20, 4000]$\,mm (the VL53L5CX
  operating range from datasheet).
\item Add per-pixel Gaussian noise with $\sigma = 5$\,mm and an
  independent 5\% per-zone dropout to the saturated max-range value (this
  models SPAD over-saturation).
\item Re-clip to $[20, 4000]$\,mm so noise cannot push values outside the
  achievable hardware range.
\end{enumerate}

The output is a $[N, 8, 8]$ tensor in millimetres. The
``MJCF-enumeration-order'' invariant (rather than alphabetical) is
load-bearing: the sensor token slot the encoder consumes must match
across simulation, training, and evaluation, and the build pipeline
controls only the source-file order.

\subsection{Data collection}
\label{sec:system-data}

We collect $N_\mathrm{traj} = 1000$ near-contact trajectories using the
TAMP planner from MolmoBot. Each rollout writes one HDF5 file with the
schema in Table~\ref{tab:hdf5-schema}. Schema-conformance and a
\emph{proximity-informative} coverage check are required before training
begins: a trajectory is proximity-informative if any sensor reads below
200\,mm at any timestep, and we require $\geq 30\%$ of trajectories to
satisfy this. The threshold was calibrated such that a standard
open-workspace PnP rollout fails it (so the data design does not
trivialise the headline comparison).

\begin{table}[t]
\centering
\caption{HDF5 episode schema.}
\label{tab:hdf5-schema}
\small
\begin{tabular}{lcll}
\toprule
key & dtype & shape & description \\
\midrule
\verb|obs/tof|   & float32 & $[T, N, 8, 8]$       & ToF in mm, clipped $[20, 4000]$ \\
\verb|obs/rgb|   & uint8   & $[T, 3, 224, 224]$    & RGB camera frame \\
\verb|obs/qpos|  & float32 & $[T, 7]$              & FR3 joint positions \\
\verb|actions|   & float32 & $[T, 7]$              & joint deltas \\
\verb|attrs.success| & bool & scalar & episode success flag \\
\verb|attrs.n_sensors| & int & scalar & sensor count, must match dataset \\
\verb|attrs.language| & str & scalar & language instruction (optional) \\
\bottomrule
\end{tabular}
\end{table}

\subsection{Normalisation}

Per-channel mean and standard deviation are computed for \verb|tof|,
\verb|qpos|, and \verb|actions| over the \emph{training} files only,
using the same seed-and-fraction split as the DataLoader. Stats are
cached to \verb|stats.json| and applied at load time; predicted action
chunks are unnormalised at inference before reaching the controller. The
unnormalisation step is non-negotiable: omitting it produces a robot
that moves at $1/\sigma$ of intended scale.

% =============================================================================
\section{Method}
\label{sec:method}

PLA fuses three streams --- proximity, vision-language, and
proprioception --- into the context of an ACT decoder. The full graph is
shown in Figure~\ref{fig:overview}. We describe each component below.

\begin{figure}[t]
\centering
\begin{tikzpicture}[
    node distance=0.8cm and 0.7cm,
    every node/.style={font=\small},
    box/.style={
        draw, rounded corners=2pt, minimum height=0.9cm, minimum width=2.6cm,
        align=center, inner sep=2pt
    },
    big/.style={
        box, minimum width=3.6cm, minimum height=1.4cm, fill=blue!7
    },
    frozen/.style={
        box, fill=gray!10, minimum width=3.6cm, minimum height=1.4cm
    },
    ctx/.style={
        box, fill=orange!10, minimum width=4.2cm
    },
    out/.style={
        box, fill=green!10, minimum width=2.6cm
    },
    arrow/.style={-{Latex[length=2mm]}, thick},
]
% inputs
\node[box] (rgb)  {RGB $[B, 2, 3, 224, 224]$};
\node[box, below=of rgb] (lang) {language str};
\node[box, below=of lang] (tof) {ToF $[B, 32, 8, 8]$ (mm)};
\node[box, below=of tof] (qpos) {qpos $[B, 7]$};
% encoders
\node[frozen, right=of rgb, yshift=-0.4cm] (vlm) {Frozen Molmo2-4B\\ + Linear};
\node[big, right=of tof] (penc) {ProximityEncoder\\ shared MLP + LN};
\node[box, right=of qpos] (pp) {Linear};
% fusion
\node[ctx, right=2.0cm of vlm] (fuse) {ModalityFusion\\ concat + LN};
% ACT
\node[big, right=of fuse, yshift=-0.7cm, fill=purple!7] (act) {ACT decoder\\ CVAE + Transformer};
% outputs
\node[out, right=of act] (out) {pred $[B, 100, 7]$};
% wires
\draw[arrow] (rgb) -- (vlm);
\draw[arrow] (lang) -- (vlm);
\draw[arrow] (tof) -- (penc);
\draw[arrow] (qpos) -- (pp);
\draw[arrow] (vlm) -- node[above, font=\scriptsize]{$[B, N_v, 512]$} (fuse);
\draw[arrow] (penc) -- node[above, font=\scriptsize]{$[B, 32, 512]$} (fuse);
\draw[arrow] (pp) -| (fuse);
\draw[arrow] (fuse) -- node[above, font=\scriptsize]{$[B, N_\mathrm{ctx}, 512]$} (act);
\draw[arrow] (act) -- (out);
\end{tikzpicture}
\caption{PLA architecture. Frozen Molmo2-4B (gray) supplies vision-language tokens;
the proximity encoder (blue) supplies one token per sensor; a proprioceptive
projection supplies one token. All three are concatenated and layer-normed,
then decoded by a 7-layer ACT transformer that predicts a 100-step joint-delta
chunk in a single forward pass. The VLM-only ACT baseline is identical except
the proximity branch is removed.}
\label{fig:overview}
\end{figure}

\subsection{Proximity encoder}
\label{sec:method-encoder}

Given $\mathbf{T} \in \R^{B \times N \times 8 \times 8}$ (normalised
proximity readings; $N = 32$), we apply a \emph{shared} two-layer MLP
followed by LayerNorm to each sensor's flattened patch:
\begin{equation}
\mathbf{p}_i = \mathrm{LN}\big(W_2\, \mathrm{ReLU}(W_1\, \mathrm{vec}(\mathbf{T}_i)) \big)
\quad\in\R^{d}
\qquad i \in \{1, \dots, N\}
\end{equation}
where $W_1 \in \R^{h \times 64}, W_2 \in \R^{d \times h}$, $h = 128$,
$d = 512$. The MLP is shared across all sensors: this matches the
inductive bias that the same hardware should apply the same encoding,
and yields $N\times$ the gradient signal per training step relative to
per-sensor encoders. Per-sensor identity is recovered downstream by the
per-token slot in the decoder context.

The output LayerNorm is essential: without it, proximity tokens have
markedly smaller magnitudes than visual-language tokens, attention
softmax saturates at the fusion boundary, and the decoder learns to
ignore proximity entirely. We have observed gradient-norm collapse
($\norm{\nabla_{W_1}}_2 < 10^{-12}$) in early prototypes lacking this
LayerNorm; the Day-6 grad-norm sanity check
(Section~\ref{sec:protocol-checks}) is calibrated to fail loud when this
recurs.

\subsection{Frozen vision-language backbone}
\label{sec:method-vlm}

Molmo2-4B~\citep{molmo2024} is loaded with all parameters frozen. Two
RGB frames (the current frame and a one-step lag) are processed through
the vision tower, mean-pooled across frames at the token level, and
projected by a trainable $\mathrm{Linear}(d_\mathrm{vlm} \to d)$ to the
shared embedding dimension. Freezing the backbone is essential for two
reasons:
\begin{enumerate}[leftmargin=1.5em,topsep=0.2em,itemsep=0.2em]
\item \emph{Catastrophic forgetting.} Fine-tuning a 4B-parameter VLM on
  $\sim$1k trajectories degrades open-vocabulary visual priors more than
  it improves manipulation accuracy.
\item \emph{Memory.} Frozen Molmo2 fits in $\sim$7\,GB; full fine-tune
  adds $\sim$16\,GB of optimiser state. The frozen configuration leaves
  room for batch size 8 on a single A100.
\end{enumerate}

\subsection{Modality fusion}
\label{sec:method-fusion}

We fuse via concatenation along the sequence dimension followed by
LayerNorm:
\begin{equation}
\mathbf{C} = \mathrm{LN}\big(\;[\mathbf{P} \,\Vert\, \mathbf{V} \,\Vert\, \mathbf{q}]\;\big)
\in \R^{B \times (N + N_v + 1) \times d}
\end{equation}
where $\mathbf{P}$ are proximity tokens, $\mathbf{V}$ are projected
visual-language tokens, and $\mathbf{q}$ is the projected proprioceptive
token. Concatenation makes no commitment about how the streams interact;
the decoder's attention layers are free to mix any subset. We
investigate a cross-attention variant (proximity tokens query
visual-language) as an ablation in Section~\ref{sec:protocol-ablations}.

\subsection{ACT decoder}
\label{sec:method-act}

We use the Action Chunking
Transformer~\citep{zhao2023learning} verbatim with the following
hyperparameters: $d=512$, 8 heads, 4 encoder layers (CVAE), 7 decoder
layers, $d_\mathrm{ffn}=3200$, latent $z \in \R^{32}$, dropout 0.1,
chunk size 100. The training loss is L1 on the action chunk plus a
constant-weighted KL term:
\begin{equation}
\mathcal{L} = \E\big[\,\norm{\hat{\mathbf{a}} - \mathbf{a}}_1\,\big]
            + \beta\,\KL\!\big(\mathcal{N}(\mu, \mathrm{diag}(e^{\log \sigma^2})) \,\Vert\, \mathcal{N}(0, I)\big),
\quad \beta = 10
\label{eq:loss}
\end{equation}
$\beta$ is held constant rather than annealed; this drives the latent
to collapse near the prior mean during training so that inference at
$z = 0$ matches the training distribution.

\subsection{The one-flag baseline}
\label{sec:method-baseline}

The VLM-only ACT baseline is constructed by setting \verb|vlm_only=True|
in the same model class. This removes the proximity encoder entirely
and the fusion produces $\mathbf{C} = \mathrm{LN}([\mathbf{V} \,\Vert\,
\mathbf{q}])$. Every other architectural detail, hyperparameter, dataset,
and seed schedule is identical between PLA and the baseline. The two
configurations differ by one line of YAML
(\verb|vlm_only: true| vs.\ \verb|false|) and we have verified at the
parameter-count level that this is the \emph{only} difference: PLA has
exactly one more module --- the proximity encoder --- whose parameter
count is removed in the baseline.

% =============================================================================
\section{Experimental Protocol}
\label{sec:protocol}

This section is locked before data collection begins. Any deviation
between this specification and the eventual reported results will be
explicitly noted in the relevant subsection of
Section~\ref{sec:results} with reasons.

\subsection{Tasks}
\label{sec:protocol-tasks}

We evaluate on four MolmoSpaces tasks, each randomised over
procthor-objaverse scenes:

\begin{enumerate}[leftmargin=1.5em,topsep=0.2em,itemsep=0.2em]
\item \textbf{pnp} (open workspace pick-and-place). Baseline competence
  test. We expect a small or zero PLA advantage here.
\item \textbf{near\_contact} (\emph{primary}). Pick-and-place with a
  rigid obstacle 5--8\,cm from the expert arm path. The obstacle is
  pre-placed in scene generation. We expect the largest PLA advantage.
\item \textbf{pnp\_color}. Multiple coloured distractors; the language
  instruction selects one. Stresses the VLM's pointing prior.
\item \textbf{pnp\_next\_to}. Spatial-relation placement (place object
  next to a reference). The hardest language test in the suite.
\end{enumerate}

\subsection{Sample sizes and pairing}
\label{sec:protocol-N}

Final per-(method $\times$ task) cell: $n = 100$ episodes. The same
seed schedule \verb|seed_base + i| for $i \in \{0, \dots, 99\}$,
\verb|seed_base = 42|, is used across all methods. Identical seeds in
MolmoSpaces produce identical scenes, object placements, and language
prompts; per-episode outcomes are therefore aligned across methods so
the paired bootstrap is valid.

\subsection{Statistical procedure}
\label{sec:protocol-stats}

For each method we compute a 95\% bootstrap confidence interval on
success rate using $B = 10\,000$ resamples with seed $0$. For each
non-baseline method we additionally compute a paired-bootstrap
two-sided p-value against the VLM-only baseline by:
\begin{enumerate}[leftmargin=1.5em,topsep=0.2em,itemsep=0.2em]
\item Forming the per-episode difference $d_i = a_i - b_i \in \{-1, 0, +1\}$
  where $a, b$ are method and baseline success vectors.
\item Recentring $\tilde d_i = d_i - \bar d$.
\item Drawing $B$ resamples of size $n$ from $\tilde d$ with replacement.
\item Reporting $p = \frac{1}{B} \sum_{b=1}^{B} \mathbf{1}\!\big[ |\bar{\tilde d}^{(b)}| \geq |\bar d| \big]$.
\end{enumerate}

\paragraph{Acceptance criterion (headline).} The headline contribution
holds if on the \texttt{near\_contact} task PLA's mean success rate
exceeds VLM-only ACT's by $\geq 5$ percentage points \emph{and} the
paired-bootstrap p-value is below $\alpha = 0.05$.

\paragraph{Multiple comparisons.} The headline test is a single
hypothesis. Secondary comparisons (4 ablations $\times$ 1 task; PLA vs
baseline $\times$ 4 tasks) are reported with Bonferroni-corrected
p-values ($\times 4$). Sensor-importance per-sensor is descriptive and
not corrected.

\subsection{Ablations}
\label{sec:protocol-ablations}

All ablations differ from PLA by exactly one YAML flag:

\begin{enumerate}[leftmargin=1.5em,topsep=0.2em,itemsep=0.2em]
\item \textbf{wrist-only}. Sensor-mask zeros out the 24 non-wrist
  sensors. Tests whether whole-body coverage is necessary or wrist
  proximity is sufficient.
\item \textbf{handcrafted}. Replace the learned shared MLP with
  hand-engineered features (per-sensor min, mean, contact-flag
  $\mathbf{1}[\min < 100\,\mathrm{mm}]$). Tests whether learning beats
  fixed statistics.
\item \textbf{conv2d}. Replace the flat MLP with a tiny per-sensor
  Conv2D over the 8$\times$8 grid. Tests whether spatial structure
  inside the SPAD grid carries useful information beyond the flat MLP.
\item \textbf{cross\_attn}. Replace concat fusion with cross-attention
  (proximity queries vision-language). Tests the "given what I see,
  which proximity readings matter" inductive bias.
\end{enumerate}

\subsection{Sensor importance}

Post-hoc and without retraining: for each sensor index $i$, mask its
input grid (zero its 8$\times$8 channel) and re-evaluate on the same
paired seeds. Importance is $\Delta_i = \mathrm{SR}_\mathrm{base} -
\mathrm{SR}_{\mathrm{masked}\,i}$. We use $n = 50$ episodes per sensor
($1600$ total).

\subsection{Pre-training sanity checks}
\label{sec:protocol-checks}

The following are pre-conditions for a valid PLA training run; failures
abort training before compute is committed:
\begin{itemize}[leftmargin=1.5em,topsep=0.2em,itemsep=0.2em]
\item \textbf{Schema and coverage:} \verb|pla.data.verify --strict|
  reports 0 NaN, schema 100\%, $\geq 30\%$ proximity-informative
  trajectories.
\item \textbf{Forward pass:} \verb|pla.checks.forward_pass| asserts
  output shapes $[B, 100, 7]$, $[B, 32]$, $[B, 32]$ for
  $\hat{\mathbf{a}}, \mu, \log\sigma^2$ on every architectural variant.
\item \textbf{Gradient flow:} the proximity-encoder L2 gradient norm
  exceeds $10^{-8}$ after 50 synthetic-batch optimisation steps. We log
  the per-step gradient norm during training; a value below this
  threshold late in training raises an explicit warning.
\end{itemize}

% =============================================================================
\section{Results}
\label{sec:results}

This section reports two classes of results: (i) Day-1 system-validation
results that are complete and verified, and (ii) the headline empirical
comparison, which is pre-registered above and \emph{will be reported on
Day 14 of the execution timeline} (see
Appendix~\ref{app:timeline}).

\subsection{System validation (Day 1)}
\label{sec:results-validation}

We have verified the following on Day 1 (synthetic-data fixtures and
\verb|DummyVLBackbone|; full transcripts in supplementary
\verb|docs/SANITY_CHECKS.md|):

\begin{itemize}[leftmargin=1.5em,topsep=0.2em,itemsep=0.2em]
\item All 49 Python source files parse cleanly; all 29 modules import
  without circular-dependency errors.
\item Forward$+$backward$+$inference passes for all 5 architectural
  variants (PLA shared-MLP, VLM-only baseline, handcrafted encoder,
  conv2d encoder, cross-attention fusion) with documented output shapes
  $[2, 100, 7]$, $[2, 32]$, $[2, 32]$.
\item Proximity-encoder gradient norm exceeds $10^{-2}$ after 20
  optimisation steps for every PLA variant
  (Table~\ref{tab:gradnorm}); the VLM-only baseline correctly reports
  ``no proximity encoder to check'' and the encoder's parameters are
  absent from the optimiser state (4{,}863{,}751 vs.\ 4{,}880{,}455
  trainable parameters at $d=64$).
\item The synthetic data pipeline is end-to-end consistent:
  6 episodes $\to$ schema validator (6/6 pass) $\to$ normaliser
  ($\mathrm{stats.json}$ written from 5 train files) $\to$ DataLoader
  (745 train / 149 val sliding windows) $\to$ 3-step training (loss
  $2.78 \to 1.55$, encoder gradient $> 0.09$).
\item Statistical estimators behave correctly on synthetic 65/50 success
  rates: bootstrap 95\% CIs $[60\%, 78\%]$ and $[33\%, 53\%]$;
  paired-bootstrap $p = 3 \times 10^{-4}$, in agreement with an
  independent $\chi^2$ check.
\item The wrist-only sensor mask zeros input indices $[8, 32)$ and
  preserves $[0, 8)$ element-wise.
\end{itemize}

\begin{table}[t]
\centering
\caption{Proximity-encoder per-parameter L2 gradient norms after 20
synthetic-batch optimisation steps. Threshold for ``learning'' is
$10^{-8}$; all observed values exceed this by 5--7 orders of magnitude.
The VLM-only baseline correctly has no proximity encoder.}
\label{tab:gradnorm}
\small
\begin{tabular}{lcccccc}
\toprule
& \multicolumn{2}{c}{shared-MLP} & \multicolumn{2}{c}{conv2d} & \multicolumn{1}{c}{handcrafted} & vlm-only \\
\cmidrule(lr){2-3}\cmidrule(lr){4-5}\cmidrule(lr){6-6}\cmidrule(lr){7-7}
parameter & PLA & cross-attn & conv2d & & handcrafted & baseline \\
\midrule
\verb|mlp.0.weight|       & $1.6\!\times\!10^{-1}$ & $2.1\!\times\!10^{-1}$ & --- & & --- & --- \\
\verb|mlp.0.bias|         & $2.5\!\times\!10^{-2}$ & $3.3\!\times\!10^{-2}$ & --- & & --- & --- \\
\verb|mlp.2.weight|       & $2.9\!\times\!10^{-1}$ & $3.5\!\times\!10^{-1}$ & --- & & --- & --- \\
\verb|conv.0.weight|      & --- & --- & $5.0\!\times\!10^{-2}$ & & --- & --- \\
\verb|conv.2.weight|      & --- & --- & $2.2\!\times\!10^{-1}$ & & --- & --- \\
\verb|proj.weight|        & --- & --- & $1.8\!\times\!10^{-1}$ & & $6.3\!\times\!10^{-3}$ & --- \\
\verb|norm.weight| (LN)   & $1.4\!\times\!10^{-2}$ & $2.0\!\times\!10^{-2}$ & $2.0\!\times\!10^{-2}$ & & $3.8\!\times\!10^{-4}$ & --- \\
\midrule
verdict & PASS & PASS & PASS & & PASS & SKIP (no encoder) \\
\bottomrule
\end{tabular}
\end{table}

\subsection{Pre-registered empirical results (Day 14)}

The result table below is the structure that will be filled. Columns
and methods are locked. The acceptance criterion
(Section~\ref{sec:protocol-stats}) is locked.

\begin{table}[t]
\centering
\caption{\pending{[PRE-REGISTERED, Day 14 deliverable]} Success rate (\%)
with bootstrap 95\% CIs and paired-bootstrap p-values vs.\ VLM-only ACT.
$n = 100$ episodes per cell, paired across methods by shared seed. Bold
entries indicate the headline test ($p < 0.05$ on \texttt{near\_contact}).}
\label{tab:results}
\small
\begin{tabular}{lcccc}
\toprule
& \texttt{pnp} & \texttt{near\_contact} & \texttt{pnp\_color} & \texttt{pnp\_next\_to} \\
\midrule
PropOnlyMLP    & \pending{--} & \pending{--} & \pending{--} & \pending{--} \\
VLM-only ACT   & \pending{--} & \pending{--} & \pending{--} & \pending{--} \\
\midrule
PLA + wrist only       & \pending{--} & \pending{--} & \pending{--} & \pending{--} \\
PLA + handcrafted      & \pending{--} & \pending{--} & \pending{--} & \pending{--} \\
PLA + conv2d           & \pending{--} & \pending{--} & \pending{--} & \pending{--} \\
PLA + cross-attn       & \pending{--} & \pending{--} & \pending{--} & \pending{--} \\
\midrule
\textbf{PLA (ours)}    & \pending{--} & \pending{--} & \pending{--} & \pending{--} \\
\bottomrule
\end{tabular}
\end{table}

% =============================================================================
\section{Discussion}
\label{sec:discussion}

\subsection{What we predict, and what would change our minds}

We predict (i) PLA $>$ VLM-only ACT on \texttt{near\_contact} by
$\geq 10$\,pp with $p < 0.05$, (ii) the gap on \texttt{pnp} (open
workspace, no nearby obstacles) is $< 5$\,pp, (iii) wrist-only $\approx$
PLA on \texttt{near\_contact}, suggesting that the wrist proximity is
the primary information source --- but with $> 5$\,pp degradation if any
non-wrist sensor is masked individually
(Section~\ref{sec:protocol-ablations}, sensor-importance), and (iv)
handcrafted features lag PLA by $\geq 5$\,pp, justifying the learned
encoder.

If the headline test \emph{fails} ($p \geq 0.05$ or $\Delta < 5$\,pp),
the candidate causes in priority order are: (a) data coverage falls
below the 30\% near-contact threshold despite a passing verifier
($\Rightarrow$ task design needs to push harder); (b) the encoder
gradient collapses late in training ($\Rightarrow$ LayerNorm position or
fusion bug); (c) the obstacle is too easy for RGB to handle anyway
($\Rightarrow$ tighten obstacle distance to 4\,cm); (d) the comparison
is genuinely null and proximity does not help in our regime
($\Rightarrow$ honest negative result; we will report it).

\subsection{Limitations}

Two caveats apply to the present scope:

\begin{enumerate}[leftmargin=1.5em,topsep=0.2em,itemsep=0.2em]
\item \emph{Simulation only.} All numbers are produced in MuJoCo with a
  calibrated VL53L5CX noise model. We do not claim sim-to-real transfer
  in this work; the noise model and clipping protocol are designed to
  make a real-world deployment plausible, not certain. Hardware
  validation is concurrent work.
\item \emph{Frozen vision backbone.} We do not adapt Molmo2 to the
  manipulation distribution. A LoRA-adapted backbone might extract more
  task-relevant visual features, narrowing the PLA gap. This is a
  question for future work and one we are explicitly trading off
  against catastrophic forgetting (Section~\ref{sec:method-vlm}).
\end{enumerate}

\subsection{Broader scientific positioning}

The paper makes a stronger claim than ``adding sensors helps''. It
argues that the \emph{representation} of proximity --- a distinct,
learnable, body-attached stream --- is what unlocks a class of tasks
that standard VLAs systematically fail at. The peripersonal regime is
not a tail of the distribution; it is the regime in which the most
common day-to-day manipulation occurs. By isolating the contribution to
a one-flag architectural change, we render the claim falsifiable in a
way that ``adding more data'' or ``adding more parameters'' would not.

% =============================================================================
\section{Conclusion}
\label{sec:conclusion}

We presented PLA, a vision-language-action policy that adds a learned
representation of whole-body time-of-flight proximity sensing to a
frozen-backbone, action-chunking architecture. The contribution is
isolated to a single architectural switch, the experimental protocol is
pre-registered with paired-bootstrap statistics, and the system has
been validated end-to-end on synthetic fixtures. The headline empirical
comparison will be reported on Day 14 of the execution timeline; the
methods, code, and acceptance criteria are locked here. We argue that
the peripersonal regime deserves a dedicated representational primitive,
and that whole-body ToF sensing combined with a shared-MLP encoder is a
strong, simple candidate for that primitive.

% =============================================================================
\section*{Reproducibility statement}

The full implementation (all model code, configs, data pipeline, eval
harness, and statistical helpers) is available at the anonymised
repository
URL\,\footnote{Anonymised for review. The release is a single Python
package importable as \texttt{pla} with one flat YAML config per
experimental condition.}.
The training-time and evaluation-time random seeds are pinned in the
configs; \texttt{stats.json} is committed alongside each checkpoint so
the exact normalisation constants are recoverable. Day-14 results will
be accompanied by per-method, per-task JSON outcome files (one per
$n=100$ run) so that the bootstrap statistics can be independently
recomputed. The anonymised repository contains
$\sim$5{,}800 lines of Python and 12 per-folder scientific READMEs that
document the contract, sanity checks, and design rationale for each
sub-package.

% =============================================================================
\bibliographystyle{plainnat}
\bibliography{references}

\appendix
% =============================================================================
\section{Execution timeline (for the pre-registration)}
\label{app:timeline}

\textbf{Today is Day 1 (May 2, 2026).} The submission deadline is Day 27
(May 28, 2026 AoE). The pre-registered execution schedule is:
\begin{itemize}[leftmargin=1.5em,topsep=0.2em,itemsep=0.2em]
\item Day 1--2: skin design + MJCF build + sanity checks (\textbf{done as of submission}).
\item Day 3: launch 1000-trajectory near-contact collection.
\item Day 4--5: train baselines.
\item Day 6: PLA training start; verify gradient flow.
\item Day 7: 50-episode quick eval (\texttt{near\_contact}, baselines).
\item Day 8--9: launch all four ablations.
\item Day 10--11: per-sensor importance sweep.
\item Day 12--13: figures.
\item \textbf{Day 14: results lock} --- Tables~\ref{tab:results} and the
  sensor-importance heatmap are filled with the n=100 numbers.
\item Day 15--21: writing.
\item Day 22--26: revision, formatting, supplementary.
\item Day 27: submit.
\end{itemize}

% =============================================================================
\section{Architectural variants studied}
\label{app:variants}

The proximity encoder, fusion module, and full PLA model are all
controlled by config flags in a single \texttt{configs/train/*.yaml}
file. The variants used in this paper are:

\begin{table}[h]
\centering
\caption{Variants and the YAML flag(s) that produce them.}
\small
\begin{tabular}{ll}
\toprule
variant & flags from \texttt{configs/train/pla.yaml} \\
\midrule
PLA (ours)       & (none --- this is the default) \\
VLM-only ACT     & \texttt{vlm\_only: true} \\
PLA + wrist-only & \texttt{sensor\_mask: [8, 9, ..., 31]} \\
PLA + handcrafted & \texttt{encoder\_type: handcrafted} \\
PLA + conv2d     & \texttt{encoder\_type: conv2d} \\
PLA + cross-attn & \texttt{fusion\_type: cross\_attn} \\
\bottomrule
\end{tabular}
\end{table}

% =============================================================================
\section{Hyperparameters}
\label{app:hyperparams}

\begin{table}[h]
\centering
\caption{Locked hyperparameters. These do not vary across configurations.}
\small
\begin{tabular}{lc}
\toprule
hyperparameter & value \\
\midrule
hidden dim $d$               & 512 \\
attention heads              & 8 \\
ACT encoder layers (CVAE)    & 4 \\
ACT decoder layers           & 7 \\
FFN dim                      & 3200 \\
latent dim $z$               & 32 \\
chunk size                    & 100 \\
RGB history $K$              & 2 \\
sensors $N$                  & 32 \\
KL weight $\beta$            & 10 (constant) \\
optimiser                    & Adam \\
learning rate                & $1 \times 10^{-5}$ \\
batch size                   & 8 \\
gradient clip                & 1.0 \\
dropout                       & 0.1 \\
\bottomrule
\end{tabular}
\end{table}

\end{document}
